\origpage{49}{}
\noindent disjoint paths of $G$ between $Z$ and $Z'$, different from $\{P_1, \dots, P_p\}$, since $e_1$ is an edge of one of them. This contradicts our hypothesis, and hence (4) follows.

From (3) and (4) it follows that $V(H) = \emptyset$, and so $z_i = z_i'$ for $1 \le i \le k$. Hence
\begin{equation*}
V(B \cap A') = V(P_1 \cup \dots \cup P_k) = Z
\end{equation*}
and so $E(B \cap A') = \emptyset$, from (1). The theorem is therefore satisfied by a 2-vertex path $P$. $\square$

The converse of (12.1) is obvious, that if there is a path-decomposition satisfying (i), (ii) and (iii)\fix{, and such that for each $t \in V(P) \setminus \{s, s'\}$ every vertex of $W(t)$ belongs to $W(e)$ for some edge $e$ of $P$ incident with $t$}, then there is a unique set of $k$ paths as in the hypothesis.\fixnote{fix:12-converse}{As printed, the converse is false: (i)--(iii) give the unique set of $k$ paths, but not that every vertex of $B \cap A'$ is on one of them, since a graph $W(t)$ may hold a vertex that is in no $W(e)$. For example, let $G$ have vertices $x, y, u$ and one edge $xy$, let $k = 1$, and let $P$ be the path $s t s'$ with $W(s) = \{x\}$, $W(s') = \{y\}$, and $W(t)$ the edge $xy$ with its ends and the vertex $u$. Then $u$ is on no path. With the added condition the converse holds, by induction on $|V(P)|$. The sketch of the proof of (12.4) uses the converse.}

Secondly we need the following.

\textbf{12.2} \textit{Let $(T, W)$ be a linked tree-decomposition of a graph $G$, with width $\le w$. Let $f_1, f_2 \in E(T)$ be distinct, and let $T_1, T_2, T_0$ be the three components of $T \setminus \{f_1, f_2\}$ where $f_i$ has ends in $T_i$ and $T_0$ ($i = 1, 2$). Let $|W(f_1)| = |W(f_2)| = k$, and let $(P, W')$ be a path-decomposition of $G$ satisfying statements (i)-(iii) of (12.1), with $W, A, B, A', B'$ replaced by $W', W(T_1), W(T_0) \cup W(T_2), W(T_0) \cup W(T_1), W(T_2)$\fix{, and such that for each $t \in V(P)$ that is not an end of $P$, every vertex of $W'(t)$ belongs to $W'(e)$ for some edge $e$ of $P$ incident with $t$}.\fixnote{fix:12.2}{As printed, (12.2) is false. Statement (12.1)(iii) does not say that $V(W'(t)) = W'(e) \cup W'(e')$, so a graph $W'(t)$ may hold an extra vertex, and then $(T^*, W^*)$ may have width $w + 1$. For example, let $G$ have vertices $a, b, y$ and one edge $ab$. Let $T$ be the path $t_1 t_0 t'_0 t_2$ with $W(t_1) = \{a\}$, $W(t_0)$ the edge $ab$ with its ends, $W(t'_0) = \{b, y\}$ and $W(t_2) = \{b\}$, which has width $1$. Let $f_1 = t_1 t_0$ and $f_2 = t'_0 t_2$, so that $k = 1$, and let $P$ be the path $s_1 t s_2$ with $W'(s_1) = \{a\}$, $W'(s_2) = \{b\}$, and $W'(t)$ the edge $ab$ with its ends and $y$. Then $W^*(t)$ has $3$ vertices. With the added hypothesis, the proof of Claim (3) works. The path-decomposition built in the proof of (12.1) has the property, since there every vertex of $B \cap A'$ is on one of the $k$ paths.} Let $P$ have ends $s_1, s_2$, where $s_i$ is the end of $f_i$ in $V(T_i)$ ($i = 1, 2$), and otherwise let $P$ be disjoint from $T$. Let $T^*$ be the tree $T_1 \cup T_2 \cup P$, and define $W^*(t) = W(t)$ if $t \in V(T_1) \cup V(T_2)$ and $W^*(t) = W'(t)$ if $t \in V(P) \setminus \{\delm{s, s'}\fix{s_1, s_2}\}$.\fixnote{fix:12.2-ends}{In (12.2) the ends of $P$ are called $s_1, s_2$. The names $s, s'$ come from (12.1), and the proof of Claim (3) writes $t \in V(P) \setminus \{s_1, s_2\}$.} Then $(T^*, W^*)$ is a linked tree-decomposition of $G$, with width $\le w$.}

\textbf{Proof.} It is easy to see that $(T^*, W^*)$ is a tree-decomposition of $G$. \fix{To see this, let $P$ have vertices $s_1 = t_0, t_1, \dots, t_m, t_{m+1} = s_2$ in order. By (12.1)(i), $W'(s_1) = W(T_1)$, $W'(s_2) = W(T_2)$, $W'(P \setminus s_1) = W(T_0) \cup W(T_2)$ and $W'(P \setminus s_2) = W(T_0) \cup W(T_1)$. Let $y$ be a vertex or an edge of $G$. If $y$ belongs to $W(t)$ for some $t \in V(T_0)$, then $y$ belongs to $W'(P \setminus s_1)$, and so to some $W'(t_j)$ with $1 \le j \le m$ or to $W(T_2)$; hence the union of the graphs $W^*(t)$ ($t \in V(T^*)$) is $G$. Let $S_y$ be the set of $t \in V(T^*)$ such that $y$ belongs to $W^*(t)$. Since $T^*$ is a tree, it suffices to show that $S_y$ induces a connected subgraph of $T^*$. The sets $S_y \cap V(T_1)$ and $S_y \cap V(T_2)$ induce subtrees, since $(T, W)$ is a tree-decomposition; and the set of $j$ ($0 \le j \le m + 1$) such that $y$ belongs to $W'(t_j)$ is an interval $I$, since $(P, W')$ is a path-decomposition. If $S_y$ meets $V(T_1)$ and also contains some other vertex of $T^*$, then $y$ belongs to $W(T_1)$ and to $W(T_0) \cup W(T_2)$, so $y$ belongs to $W(t)$ for some $t \in V(T_1)$ and to $W(t')$ for some $t' \notin V(T_1)$, and so to $W(s_1)$, since $s_1$ is on the path of $T$ between $t$ and $t'$; and then $0 \in I$. Similarly, if $S_y$ meets $V(T_2)$ and contains some other vertex, then $s_2 \in S_y$ and $m + 1 \in I$. Since $S_y \cap \{t_1, \dots, t_m\} = \{t_j : j \in I, 1 \le j \le m\}$, it follows that $S_y$ induces a connected subgraph of $T^*$. Hence $(T^*, W^*)$ is a tree-decomposition of $G$.}\fixnote{fix:12.2-tree-proof-added}{The paper leaves this check to the reader. The argument above is the one used in the formalization.} We shall check its width, and check that it is linked. Since $(P, W')$ satisfies (12.1)(iii), it follows that

\textit{(1) There are $k$ mutually vertex-disjoint paths of $G$ between $W(f_1)$ and $W(f_2)$.} Let $Q$ be the path of $T$ with first and last edges $f_1$ and $f_2$. From (1), we have

\textit{(2) $|W(e)| \ge k$ for each edge $e$ of $Q$.}

\textit{(3) $(T^*, W^*)$ has width $\le w$.}

\textit{Subproof.} If $t \in V(T_1) \cup V(T_2)$ then $|V(W^*(t))| = |V(W(t))| \le w + 1$, and if $t \in V(P) \setminus \{s_1, s_2\}$, then\fix{, since every vertex of $W'(t)$ is in $W'(e) \cup W'(e')$ by hypothesis, where $e, e'$ are the edges of $P$ incident with $t$, it follows from (12.1)(iii) that}
\begin{equation*}
|V(W^*\fix{(t}))| = |V(W'(t))| \le k + 1 = |W(f_1)| + 1 \le w + 2.\fixmark{fix:12.2-claim3}
\end{equation*}
\fixnotetext{The bound $|V(W'(t))| \le k + 1$ needs the hypothesis added to (12.2): in case (a) of (12.1)(iii), $V(W'(t)) = W'(e)$ has $k$ vertices, and in case (b), $V(W'(t)) = X \cup \{v, v'\}$ has $k + 1$. The left side of the display is also a misprint for $|V(W^*(t))|$.}
